% Hipótesis de investigación (hipotesis.tex y documento completo).
% Sigue el material de clase «Hipótesis o supuestos» de Taller de
% Investigación (Dra. Morales Rosales): primero se decide si corresponde una
% hipótesis, un supuesto o ninguno, y después se formula con sus cinco pasos
% (pregunta, variables, antecedentes y marco teórico, respuesta tentativa y
% congruencia). Redacción en tercera persona.
% Por indicación del equipo solo se presenta la hipótesis de investigación;
% la hipótesis nula queda redactada y comentada al final de la subsección
% «Hipótesis de Investigación» para usarla cuando se defina la prueba
% estadística.
% En hipotesis.tex y en trabajo.tex los apartados son subsecciones de
% «Hipótesis», con los subtítulos alineados a la izquierda (nivel 2 de APA).
\section{Hipótesis}
\let\hipsec\subsection

\hipsec{Elección entre Hipótesis y Supuestos}
La decisión de formular una hipótesis, un supuesto o ninguno depende de la
pregunta de investigación, los objetivos, el alcance, el enfoque y el diseño
metodológico. En este proyecto la pregunta interroga por el efecto de un
sistema de validación con códigos QR dinámicos en el control de acceso
peatonal, y el objetivo general consiste en evaluar ese efecto. El enfoque es
mixto con predominancia cuantitativa, el alcance llega a ser explicativo y el
diseño es cuasiexperimental, con una medición antes del piloto y otra durante
él en el mismo acceso. Lo que se busca, por lo tanto, es contrastar con datos
una diferencia esperada entre el prototipo y el método convencional de
revisión visual de credenciales.

Por estas razones corresponde formular una hipótesis. En los estudios
cuantitativos se plantean hipótesis cuando el alcance es correlacional o
explicativo, o cuando se pronostica un dato, mientras que los estudios
exploratorios no las requieren \parencite{hernandezsampieri2014metodologia}.
No se trata de un supuesto, porque la investigación no parte de una premisa
cualitativa que solo oriente la interpretación, sino de una proposición que
puede medirse con la hoja de observación del Apéndice~C. Tampoco se trata de
un estudio que no requiera hipótesis, ya que no se limita a construir el
prototipo y comprobar que funciona. Esa parte del trabajo, que incluye el
rechazo de códigos vencidos, repetidos o alterados y el tiempo de lectura en
banco, se verifica con criterios de aceptación y no con una hipótesis de
investigación.

\hipsec{Pregunta de Investigación y Variables}
La hipótesis responde a la pregunta planteada al cierre del planteamiento
del problema:
\begin{quote}
	¿Cuál es el efecto de implementar un sistema automatizado de
validación mediante códigos QR dinámicos en la eficiencia general del control
de acceso peatonal de los estudiantes del TESCHA?

\end{quote}
La variable independiente es el mecanismo de validación de identidad, con dos
condiciones: la revisión visual de la credencial física, vigente en la línea
base, y la lectura del código QR dinámico del prototipo, vigente durante el
piloto. Las variables dependientes son las dos manifestaciones de la
eficiencia del acceso que se observan en el sitio. La primera es el tiempo
promedio de validación por estudiante, en segundos, contado desde que el
estudiante llega al punto de revisión hasta que lo cruza. La segunda es la
longitud máxima de la fila, en número de personas, registrada en cada
intervalo de cinco minutos.

El contexto es el acceso peatonal piloto del \tescha{} en los minutos previos
al inicio de clases del turno vespertino. Para no atribuir al prototipo un
cambio que provenga de otra causa, se mantienen el mismo acceso, la misma
franja horaria, los mismos observadores y el mismo instrumento, y se registra
el número de puntos de revisión activos. Las incidencias de identidad no
forman parte de la hipótesis: se describen como información complementaria
sobre la seguridad del acceso, porque un aumento de los casos detectados
puede deberse a una mayor capacidad de detección y no a un peor control.

\hipsec{Fundamentación}
% Descartada: Tan et al. (2023), DOI 10.1145/3629296.3629331. Estima 30 s
% por estudiante con registro manual y 20 s con QR, pero trata de asistencia
% en exámenes, es de 2023 y no se pudo leer el texto completo (ACM bloquea la
% descarga) para confirmar cómo se obtuvieron esos tiempos.
La evidencia más directa proviene de \textcite{john2026qrgate}, quienes
sustituyeron la verificación manual de identidad en la puerta principal del
Dar es Salaam Maritime Institute, en Tanzania, por un sistema con código QR.
En una prueba con 50 estudiantes, el tiempo medio de verificación bajó de
aproximadamente 20 segundos a menos de 5 segundos por estudiante. Los
antecedentes del proyecto apuntan en el mismo sentido.
\textcite{cano2026sistema} reportan menos tiempo dedicado al registro de
asistencia al sustituir el método tradicional por un sistema con código QR en
la educación superior, aunque sin cifras. \textcite{tinajero2026acceso}
muestran que un lector de bajo costo conectado a la base de datos
institucional atendió un flujo aproximado de 4,500 estudiantes diarios con
tiempos de respuesta menores a un segundo. Por su parte,
\textcite{alsawi2026qracs} concluyen que los sistemas de control de acceso
con código QR pueden mejorar la eficiencia operativa. En el marco teórico se
expone, además, que la revisión visual depende de la atención del personal de
vigilancia \parencite{torres2019qr} y que el tiempo de espera de los
estudiantes puede optimizarse en los servicios escolares
\parencite{talavera2025espera}.

La relación entre las dos variables dependientes se desprende del análisis
del Apéndice~C. La utilización del acceso es $\rho = \lambda / (c\,\mu)$, donde
la tasa de servicio $\mu$ es inversa al tiempo medio de validación. Si este
tiempo disminuye con la misma llegada de estudiantes ($\lambda$) y el mismo
número de puntos de revisión ($c$), la utilización baja y la fila tiende a
acortarse. Por eso se espera que ambas variables cambien en el mismo sentido.

Los mismos estudios muestran los límites de esta expectativa. La reducción
reportada por \textcite{john2026qrgate} se obtuvo en un entorno controlado,
en una sola puerta y con una muestra pequeña. Los autores advierten que no
refleja necesariamente el funcionamiento con más carga y que las fallas de
conexión llegaron a retrasar la verificación. Además, ninguno de los estudios
midió la longitud de la fila. En el estudio más cercano al lector del
prototipo, el 30\,\% de las validaciones falló por iluminación insuficiente
o por la posición de la cámara \parencite{henriana2026gate}. Si abrir el código en el
teléfono o repetir una lectura fallida toma más tiempo que mostrar la
credencial, los resultados del piloto no respaldarán la hipótesis, y ese
resultado también responde a la pregunta de investigación.

\hipsec{Hipótesis de Investigación}
Con base en lo anterior, se formula la siguiente hipótesis de investigación:

\medskip
{\setlength{\fboxsep}{8pt}%
\noindent\fbox{\parbox{\dimexpr\linewidth-2\fboxsep-2\fboxrule\relax}{%
		\justifying
		La implementación de un prototipo de validación con códigos QR dinámicos
		reduce el tiempo promedio de validación por estudiante y la longitud
		máxima de la fila en el acceso peatonal piloto del \tescha{} durante el
		turno vespertino, en comparación con la revisión visual de credenciales.}}}
\medskip

La hipótesis expresa una diferencia entre dos condiciones del mismo acceso y
no una relación entre variables medidas en una sola ocasión. Se contrastará
al comparar la línea base con el piloto mediante la hoja de observación del
Apéndice~C.

% Hipótesis nula, redactada para cuando se defina la prueba estadística.
% La hipótesis nula niega o refuta la relación que propone la hipótesis de
% investigación (Hernández-Sampieri et al., 2014, cap. 6). Para activarla,
% quitar los signos de comentario de las líneas siguientes.
%
% Para el contraste estadístico, la hipótesis nula correspondiente es la
% siguiente:
%
% \medskip
% {\setlength{\fboxsep}{8pt}%
% \noindent\fbox{\parbox{\dimexpr\linewidth-2\fboxsep-2\fboxrule\relax}{%
% 		\justifying
% 		La implementación de un prototipo de validación con códigos QR dinámicos
% 		no reduce el tiempo promedio de validación por estudiante ni la
% 		longitud máxima de la fila en el acceso peatonal piloto del \tescha{}
% 		durante el turno vespertino, en comparación con la revisión visual de
% 		credenciales.}}}
% \medskip
%
% La hipótesis nula se rechazará solo si la diferencia observada entre la
% línea base y el piloto es estadísticamente significativa
% \parencite{hernandezsampieri2014metodologia}.

\hipsec{Congruencia con los Objetivos y la Metodología}
La hipótesis corresponde a la pregunta, porque responde de forma tentativa
cuál será el efecto del sistema, y la eficiencia del acceso se expresa en las
dos variables definidas como indicadores principales. También
corresponde al objetivo general, que consiste en evaluar ese efecto, y a los
objetivos específicos primero y último: el primero establece la línea base
del tiempo de validación y de la fila, y el último compara esa línea base con
el piloto. Por último, es congruente con el diseño cuasiexperimental de
preprueba y posprueba con un solo grupo, que permite observar la diferencia
pero no descarta por completo otras causas, como un cambio en la afluencia.
Por esa razón, si los datos respaldan la hipótesis, el resultado se
interpretará como evidencia orientativa para el caso estudiado y no como una
generalización.
